\clearpage
\section{Validation}

The implementation is covered by unit and integration tests.
\begin{itemize}
    \item \texttt{ControlUnitActorTest} checks the state machine, arming modes, timing behavior, and recovery rules.
    \item \texttt{SensorActorTest} and \texttt{KeypadActorTest} check receptionist discovery and message delivery.
    \item \texttt{ClusteredSystemTest} verifies three-node cluster formation, sensor handling, and restart recovery.
    \item \texttt{NodeStartupTest} checks the startup parser and cluster seed-node generation.
\end{itemize}

The tests are useful not only because they confirm correctness, but because
they document the intended distributed behavior. In particular, they lock in
the three-node setup, the recovery semantics, and the fact that the control
unit never assumes a dangerous default after a restart.

\subsection{Manual Multi-Node Execution Example}

The system can also be validated manually by starting three separate cluster
nodes on localhost and then driving one distributed scenario end-to-end.

\paragraph{Startup commands}
Open three terminals in the module directory and run:
\begin{verbatim}
./run-cshas.sh control-unit --host 127.0.0.1 --port 2551 \
  --seed-nodes 127.0.0.1:2551,127.0.0.1:2552,127.0.0.1:2553

./run-cshas.sh keypad --host 127.0.0.1 --port 2552 \
  --seed-nodes 127.0.0.1:2551,127.0.0.1:2552,127.0.0.1:2553

./run-cshas.sh sensor --host 127.0.0.1 --port 2553 \
  --sensor-id front_door --sensor-type DOOR_WINDOW --zone PERIMETER \
  --seed-nodes 127.0.0.1:2551,127.0.0.1:2552,127.0.0.1:2553
\end{verbatim}

On Windows PowerShell, the same scenario can be started with
\texttt{.\textbackslash run-cshas.ps1} and the same arguments.

\paragraph{Interactive scenario}
After the three nodes join the cluster, the following commands are sufficient
to observe the main distributed transitions:
\begin{enumerate}
    \item On the control-unit terminal, type \texttt{status}. The expected state is \texttt{STARTUP\_RECOVERY}.
    \item On the keypad terminal, type \texttt{pin 1234}. This leaves the safe recovery mode.
    \item On the control-unit terminal, type \texttt{status}. The expected state is now \texttt{DISARMED}.
    \item On the keypad terminal, type \texttt{arm full 1234}. This starts the exit delay.
    \item On the control-unit terminal, type \texttt{status}. The expected state is \texttt{EXIT\_DELAY}, and shortly after that it becomes \texttt{ARMED}.
    \item On the sensor terminal, press ENTER. This simulates a distributed sensor activation.
    \item On the control-unit terminal, type \texttt{status}. The expected state is \texttt{ENTRY\_DELAY}.
    \item On the keypad terminal, type \texttt{pin 1234} before the entry timeout expires.
    \item On the control-unit terminal, type \texttt{status}. The expected state returns to \texttt{DISARMED}.
\end{enumerate}

\paragraph{Expected observations}
This experiment shows that:
\begin{itemize}
    \item the cluster is formed by three separate node processes;
    \item keypad and sensor nodes discover the control unit remotely;
    \item arming and disarming commands travel across nodes only through actor messages;
    \item a remote sensor activation changes the control-unit state machine exactly as in the non-clustered assignment.
\end{itemize}
